% !TeX root = 16_002_Fall24.tex
\def\Lect{I}
\usepackage[OT2,T1]{fontenc}
\usepackage{setspace}
\onehalfspacing

\usepackage{graphicx,color,multirow,subcaption,wrapfig,amsmath,amssymb,amsthm}
\usepackage[margin=1in]{geometry}
\usepackage[english]{babel}
\usepackage[latin1]{inputenc}
\usepackage{lmodern,array,stackrel,esint,textgreek}

\usepackage[pdftex, bookmarksopen=true, colorlinks=true, linkcolor=blue, citecolor=blue, urlcolor=blue,filecolor=blue, pagebackref=false, hypertexnames=false, plainpages=false, pdfpagelabels ]{hyperref}

\usepackage{etex,setspace,pbox}
\let\checkmark\relax 
\usepackage{dingbat,wasysym,ifpdf,pdfpages,mathtools,xcolor}
\usepackage{listings}
\lstset{ 
    basicstyle=\ttfamily\scriptsize,    % The basic font style
    language=Python,                     % Default language for the Python listings below
    numbers=left,                        % Show line numbers
    numberstyle=\tiny\color{gray},      % Style for line numbers
    keywordstyle=\color{blue},            % Keywords color
    stringstyle=\color{magenta},              % Strings color
    numbersep=10pt,                       % Distance between line numbers and code
    backgroundcolor=\color{lightgray!10}, % Background color for the code block
    frame=single,                         % Frame around the code
    showstringspaces=false,               % Don't mark spaces inside strings
    breakatwhitespace=true,               % Break lines only at whitespace
    captionpos=b,                         % Caption position (b for below, t for above)
    escapeinside={\%*}{*)},               % Allows escaping to LaTeX within code
}


\usepackage[thicklines]{cancel}
\usepackage[mmddyy,nodayofweek,12hr]{datetime}
\usepackage{algorithm,algorithmic,blkarray}
\usepackage{textcomp,latexsym,amssymb,aircraftshapes}
\usepackage{pagecolor}% http://ctan.org/pkg/{pagecolor,lipsum}
\usepackage[percent]{overpic}
\usepackage{cleveref}

% removes beamer captions
\usepackage{caption}
%\captionsetup{labelformat=empty,labelsep=none}
%\renewcommand{\figurename}{} 

\usepackage[round, comma]{natbib}       % use Natbib for author-name
                                        % citations. In a presentation,
                                        % this is more legible than pure
                                        % numbers, like [21]. Google
                                        % "natbib" for more information.
                                        % Instead of "\cite{abc}", use
                                        % "\citep{abc}" or "\citet{abc}"
\bibliographystyle{abbrvnat}            % Natbib specific
%\bibliographystyle{ieeetr}            % Natbib specific
%\bibpunct{(}{)}{,}{a}{}{,}              % Natbib specific
\bibpunct{[}{]}{,}{a}{}{,}              % Natbib specific
\setlength{\bibhang}{10pt}
\setlength{\bibsep}{1pt}%

\makeatletter
\DeclareRobustCommand\citep
{\begingroup\tiny\NAT@swatrue\let\NAT@ctype\z@\NAT@partrue
	\@ifstar{\NAT@fulltrue\NAT@citetp}{\NAT@fullfalse\NAT@citetp}}
\makeatother

\newif\ifPDF
\ifx\pdfoutput\undefined\PDFfalse \else\ifnum\pdfoutput > 0 \PDFtrue
\else\PDFfalse \fi \fi

\usepackage{tikz,pgfplots}
\usetikzlibrary{calc,fadings,decorations.pathreplacing,quotes,angles,shapes,positioning,automata,arrows,fit,backgrounds,calc,shapes.misc,trees}
\usetikzlibrary{decorations.pathmorphing,decorations.text,decorations.markings,patterns}
%\tikzstyle{block} = [draw, fill=blue!20, rectangle,minimum height=1em, minimum width=2em]
\tikzstyle{sum} = [draw, fill=blue!20, circle, node distance=1cm]
\tikzstyle{input} = [coordinate]
\tikzstyle{output} = [coordinate]
\tikzstyle{pinstyle} = [pin edge={to-,thin,black}]
\tikzstyle{block} = [draw, fill=blue!20, rectangle,  minimum height=3em, minimum width=6em]
\tikzstyle{sblock} = [draw, fill=blue!20, rectangle,  minimum height=1em, minimum width=2em]
\tikzstyle{int}=[draw, fill=blue!20, minimum size=5em]
\tikzstyle{init} = [pin edge={to-,thin,black}]

\usetikzlibrary{decorations.pathmorphing,patterns} % springs
\usetikzlibrary{spy}

\usetikzlibrary{circuits.ee.IEC, positioning}
\usepackage[american voltages, american currents,siunitx]{circuitikz}

\tikzset{My Arrow Style/.style={single arrow, fill=red!50, anchor=base, align=center,text width=2.8cm}}
\newcommand{\MyArrow}[2][]{\tikz[baseline] \node [My Arrow Style,#1] {#2};}

\usepackage[most]{tcolorbox}
\tcbset{fonttitle=\sffamily\bfseries\Large}
\tcbuselibrary{theorems}
\tcbset{
defstyle/.style={fonttitle=\bfseries\upshape\scriptsize, fontupper=\slshape\footnotesize, fontlower=\footnotesize,arc=0mm, colback=blue!0,colframe=blue!55!red},
theostyle/.style={fonttitle=\bfseries\upshape, fontupper=\slshape\footnotesize,fontlower=\footnotesize,colback=blue!10,colframe=blue!85!black},
}
%\tcbmaketheorem{defin}{Definition}{defstyle}{mytheorem}{def}
%\tcbmaketheorem{thm}{Theorem}{theostyle}{mytheorem}{theo}

\usepackage{titlesec}
% spacing: how to read {12pt plus 4pt minus 2pt}
%           12pt is what we would like the spacing to be
%           plus 4pt means that TeX can stretch it by at most 4pt
%           minus 2pt means that TeX can shrink it by at most 2pt
%       This is one example of the concept of, 'glue', in TeX
%\titlespacing{command}{left spacing}{before spacing}{after spacing}[right]
\titlespacing\section{0pt}{0pt plus 4pt minus 2pt}{-2pt plus 2pt minus 2pt}
\titlespacing\subsection{0pt}{6pt plus 4pt minus 2pt}{-6pt plus 2pt minus 2pt}
\titlespacing\subsubsection{0pt}{6pt plus 4pt minus 2pt}{-12pt plus 2pt minus 2pt}

\usepackage{tocloft}
\setlength\cftparskip{-2pt}
\setlength\cftbeforesecskip{1pt}
\setlength\cftaftertoctitleskip{2pt}

\usepackage{empheq}
%\usepackage[most]{tcolorbox}
\newtcbox{\mymath}[1][]{%
	nobeforeafter, math upper, tcbox raise base,
	enhanced, colframe=blue!30!black,
	colback=blue!30, boxrule=1pt,
	#1}
	
\pgfplotsset{
	tick label style={font=\small},  
	legend entry/.initial=,
	every axis plot post/.code={%
		\pgfkeysgetvalue{/pgfplots/legend entry}\tempValue
		\ifx\tempValue\empty
%		\pgfkeysalso{/pgfplots/forget plot}%
		\else
		\expandafter\addlegendentry\expandafter{\tempValue}%
		\fi
	},
}

\usepackage{polynom}
\polyset{vars=s} % make X, Y, Z, and \xi into variables	

\author{Jonathan P. How\\ Massachusetts Institute of Technology\\ 
\texttt{jhow@mit.edu}}
\date{}
\newcounter{lect}
\counterwithin{figure}{lect}
\renewcommand{\thesection}{\thelect.\arabic{section}}
\usepackage{titling}
\usepackage{titletoc}
\dottedcontents{subsection}[5.5em]{}{3.2em}{1pc}
\dottedcontents{section}[1cm]{}{3em}{1em}
